\documentclass[11pt]{article}
\usepackage{amsmath,amssymb,amsthm}
\newtheorem{theorem}{Theorem}
\begin{document}
\title{An exact compact eight-wire adaptive-residue decoder}
\author{}
\date{}
\maketitle

\paragraph{Verification status.}
This is a finite, exact computational theorem, with an elementary written
argument extending the certificate to every sufficiently small positive
scale. The accompanying checker uses rational arithmetic and an explicit
list of 256 active bases. It is not a Lean proof. The coefficient height
argument and the full-cube embedding below are written proofs. No growing
capacity theorem, logarithmic saving, or bound on worstCaseHeadComplexity
is claimed. In particular this does not improve the older eight-pole,
twenty-query finite construction; its additional feature is the fixed
butterfly structure and strictly positive target-dependent residues.

\paragraph{Notation.}
There are eight wires numbered $0,\ldots,7$, whose binary bits have
indices $0,1,2$. A butterfly layer $k$ switches bit $k$, independently
for each lower prefix $s<2^k$ and higher block $g<2^{2-k}$. Let
$\theta_{k,s,g}\in\{0,1\}$ be its twelve switch settings, and let
$\pi_\theta$ be the resulting input-to-output permutation, with layers
applied in order $0,1,2$. Set
\[
 u_{k,g}(h)={\bf1}_{\lfloor h/2^{k+1}\rfloor=g}(-1)^{h_k},\qquad
 p_y(\epsilon)=\sum_{j=0}^2 y_j\epsilon^{j+1},\qquad
 t_{k,s,g}(\epsilon)=1-p_s(\epsilon).
\]
The query times may repeat. Their Haar profiles distinguish the queries.

\begin{theorem}[Positive adaptive residues]
For every $0<\epsilon\leq2^{-240}$ and every butterfly setting $\theta$,
there are strictly positive real numbers $c_0,\ldots,c_7$ such that
\[
 (1-2\theta_{k,s,g})\sum_{h=0}^7
 \frac{u_{k,g}(h)c_{\pi_\theta(h)}}
 {t_{k,s,g}(\epsilon)-p_{\pi_\theta(h)}(\epsilon)}>0
 \qquad\hbox{for all twelve queries.}                 \tag{1}
\]
The poles and times are fixed functions of $\epsilon$, independent of
$\theta$. The residues and permutation depend on the target.
\end{theorem}
\begin{proof}
Let $A_\theta(\epsilon)$ be the $12$ by $8$ signed coefficient matrix
in (1), with columns indexed by output wire $y$. Its entries are
\[
 \frac{(1-2\theta_{k,s,g})u_{k,g}(\pi_\theta^{-1}(y))}
 {1-p_s(\epsilon)-p_y(\epsilon)}.                    \tag{2}
\]
This matrix does not depend on the four layer-zero switch bits. Indeed,
flipping such a switch interchanges the two corresponding input wires
before the unchanged later routing. All higher-layer Haar profiles are
constant on this pair. The layer-zero profile of that block changes sign,
and its desired sign changes sign too. Other layer-zero profiles vanish
on the pair. Thus only the eight higher-layer bits matter.

At $\epsilon_*=2^{-240}$, the file
\texttt{butterfly-adaptive-haar-check.json} supplies one active basis for
every one of the $256$ higher-layer words. A basis selects eight rows
from the twenty inequalities
\[
 A_\theta(\epsilon_*)c\geq {\bf1},\qquad c\geq {\bf1}.
\]
The checker reconstructs the unique solution to the eight selected
active equalities. It verifies that the basis is nonsingular, every
coordinate of $c$ is positive, and every row of $A_\theta c$ is positive.
It also checks the exact matrix equality for all sixteen choices of
bottom-layer bits. This gives $49\,152$ strict signed-score checks.
The certificate contains bases, not trusted LP answers: verification
uses exact rational Gaussian elimination and does not invoke an LP solver.

Here is a uniform extension from the certificate scale to all smaller
scales. Each denominator
$D_{ey}(\epsilon)=1-p_s(\epsilon)-p_y(\epsilon)$ is an integer
polynomial whose coefficient $\ell^1$ norm is at most seven. It is
positive for $0<\epsilon\leq1/10$. Multiply query row $e$ by
$L_e=\prod_{y=0}^7D_{ey}$. The resulting coefficients have norm at most
$7^7$, and the active-equality right side has norm at most $7^8$.
Coordinate equalities have coefficient and right-side norm one.
For any certified active basis, write the cleared system as
$B(\epsilon)c=b(\epsilon)$ and use Cramer's rule
$c_y=N_y(\epsilon)/\Delta(\epsilon)$.
The determinant $\Delta$ and all eight Cramer numerators $N_y$ are
integer polynomials with norm at most $8!7^{64}$. The numerator of any
cleared score is a sum of eight products, so its norm is at most
\[
 8\cdot8!\cdot7^{71}<2^{230}.                       \tag{3}
\]
The same bound covers $\Delta$ and every $N_y$.

A nonzero integer polynomial of coefficient norm less than $K$ has no
root in $(0,1/(K+1))$. Factor its lowest nonzero power; the remaining
constant coefficient has absolute value at least one, whereas the
absolute value of its tail is at most $Kx<1$ there. Apply this observation
with $K=2^{230}$ to the determinant, all coordinate numerators, and all
score numerators. Each is nonzero, since its value at $\epsilon_*$ is
nonzero. None changes sign on $(0,\epsilon_*]$. Their certified signs
therefore give positive coordinates and strict scores throughout this
interval. This proves (1).
\end{proof}

\paragraph{Strict full-cube realization.}
For every row count $r\geq1$, the theorem yields any $r$ by $12$ sign
table with sixteen strictly admissible reverse heads and exact zeros
on both Boolean axes. Here are the budget details, to keep grid positivity
separate from full-cube positivity.
Choose
\[
 0<\epsilon<\min\{2^{-240},1/(100(r+12))\},\quad
 \alpha=1-3\epsilon,\quad
 \widetilde p_y=\alpha(p_y+\epsilon),\quad
 \widetilde t_e=\alpha(t_e+\epsilon).
\]
Then $0<\widetilde p_y<3\epsilon$,
$0<1-\widetilde t_e<4\epsilon$, and
$\widetilde t_e-\widetilde p_y=\alpha(t_e-p_y)>0$.
For each target row $i$, choose the theorem's permutation $\pi_i$ and
residues $c_{iy}>0$. Put $q_{hi}=\widetilde p_{\pi_i(h)}$ and
$a_{hi}=\alpha c_{i,\pi_i(h)}$. The query profile $u_e(h)$ is fixed,
independent of the row. For each head index $h$, pair two denominators
\[
 B_h^\pm(z,w)=1+\epsilon-\sum_i q_{hi}z_i
 -\sum_e\bigl(\epsilon+
       {\bf1}_{u_e(h)=\pm1}(1-\widetilde t_e)\bigr)w_e.
\]
Both have strictly negative slopes. The total magnitude of their slopes
is at most $3r\epsilon+12\epsilon+48\epsilon<1$, so they are positive
on the entire Boolean cube. Give the two heads opposite affine
numerators $\sum_i(1-q_{hi})a_{hi}z_i$.
The numerators vanish on $z=0$. On $w=0$ the denominators coincide,
so the pair cancels exactly. At a one-hot grid point, the paired identity
\[
 \frac{(1-q)a}{\widetilde t-q}-\frac{(1-q)a}{1-q}
 =\frac{(1-\widetilde t)a}{\widetilde t-q}
\]
shows that the pair contributes
$(1-\widetilde t_e)u_e(h)a_{hi}/(\widetilde t_e-q_{hi})$.
The positive common query factor preserves (1). This is the signed-mask
variant of the already Lean-checked pairedRowsAxisData construction in
\texttt{formalization/HeadComplexity/Results/ReverseMaskedCauchy.lean}.
The finite shattering theorem and this instantiated sixteen-head
corollary have not themselves been formalized.

\paragraph{What an induction still needs.}
The same pole hierarchy at sixteen wires has target words with an exact
Farkas obstruction at several tested scales. Faster bit scales can remove
those examples, but that is not a universal theorem. Independent positive
residue witnesses for two subtrees do not imply that their feasibility
cones intersect. Neither a positive-residue induction, a faster-hierarchy
induction, nor the version with one bias per input Haar block has been
proved. These are the specific gaps between this finite certificate and
an unbounded-capacity local compiler.
\end{document}
